% Part: first-order-logic
% Chapter: natural-deduction
% Section: rules-and-proofs

\documentclass[../../../include/open-logic-section]{subfiles}

\begin{document}

\iftag{FOL}
      {\olfileid[hi]{fol}{ntd}{rul}}
      {\olfileid[hi]{pl}{ntd}{rul}}

\olsection{नियम और \usetoken{p}{derivation}}

\begin{explain}
स्वाभाविक-निगमन तंत्रों का उद्देश्य गणितीय प्रमाण में प्रयुक्त अनौपचारिक
तर्क-प्रक्रिया के बहुत निकट चलना है (इसीलिए यह कुछ हद तक ``स्वाभाविक'' है)।
स्वाभाविक-निगमन प्रमाण मान्यताओं से आरम्भ होते हैं। फिर अनुमान-नियम लागू किए
जाते हैं। मान्यताओं को \Intro{\lnot}, \Intro{\lif}, \iftag{FOL}{\Elim{\lor}
  और \Elim{\lexists}}{ और \Elim{\lor}} अनुमान-नियमों द्वारा
``!!{discharged}'' किया जाता है, और स्पष्टता के लिए !!{discharged} मान्यता का
लेबल अनुमान के पास रखा जाता है।
\end{explain}

\begin{defn}[मान्यता]
किसी भी शाखा की सबसे ऊपर की स्थिति में स्थित कोई भी !!{sentence}
\emph{मान्यता} है।
\end{defn}

स्वाभाविक निगमन में प्रत्येक !!^{derivation} !!{sentence}s का एक विशेष वृक्ष
होती है। उसमें सबसे ऊपर के !!{sentence}s को मान्यताएँ कहा जाता है; और यदि !!a{sentence} एक,
दो या तीन अन्य अनुक्रमों के नीचे स्थित हो, तो उसे किसी अनुमान-नियम के सही
प्रयोग से निकलना चाहिए। अनुमान के ऊपर स्थित !!{sentence}s को उसके
\emph{पूर्वाधार} और नीचे स्थित !!{sentence} को अनुमान का \emph{निष्कर्ष} कहते
हैं। नियम युग्मों में आते हैं—प्रत्येक !!{operator} के लिए एक प्रवेश-नियम और
एक निरसन-नियम। वे निष्कर्ष में !!a{operator} प्रविष्ट करते हैं या नियम के किसी
पूर्वाधार से !!a{operator} हटाते हैं। कुछ नियम किसी निश्चित प्रकार की मान्यता
को \emph{!!{discharged}} करने देते हैं। कौन-सी मान्यता किस अनुमान द्वारा
!!{discharged} की जाती है, यह दिखाने के लिए हम मान्यता और अनुमान दोनों को
लेबल भी देते हैं। इसका संकेत मान्यता को ``$\Discharge{!A}{n}$.'' के रूप में
लिखकर दिया जाता है।

% Only include this sentence is on of \land, lor, lif or lnot is defined.

\iftag{notprvNot,notprvAnd,notprvOr,notprvIf}{}%
	{सभी !!{operator}s $\land$, $\lor$, $\lif$, $\lnot$, और $\lfalse$ के लिए नियमों पर विचार करना प्रचलित है, भले ही उनमें से कुछ परिभाषित हों।}



\end{document}
